\DocumentMetadata{
  pdfversion=2.0,
  pdfstandard=ua-2,
  lang=en-US,
  tagging=on,
  tagging-setup={math/setup=mathml-SE,tabsorder=structure}
}
\documentclass[11pt,letterpaper]{article}
\usepackage[margin=0.68in,includefoot]{geometry}
\usepackage{newcomputermodern}
\usepackage{amsmath,amscd,mathtools}
\usepackage{tikz,adjustbox}
\providecommand{\square}{\mdlgwhtsquare}
\usepackage[hidelinks]{hyperref}
\hypersetup{
  pdftitle={Analysis PhD Exam, August 2022},
  pdfauthor={University of Florida Department of Mathematics},
  pdfsubject={Graduate mathematics examination},
  pdfdisplaydoctitle=true,
  bookmarksopen=true
}
\setcounter{secnumdepth}{0}
\setlength{\parindent}{0pt}
\setlength{\parskip}{0.28em}
\setlength{\emergencystretch}{3em}
\setlength{\leftmargini}{2.15em}
\setlength{\leftmarginii}{2.65em}
\setlength{\leftmarginiii}{3em}
\pagestyle{empty}
\raggedbottom
\allowdisplaybreaks
\NewTaggingSocketPlug{block/list/label}{examproblem}{%
  \tagstructbegin{tag=Lbl}\tagstructend%
  \tagstructbegin{tag=\LBody}%
  \tagstructbegin{tag=\ProblemHeadingTag,title-o={\ProblemHeadingText}}%
  \tagmcbegin{tag=\ProblemHeadingTag,actualtext={\ProblemHeadingText}}%
  #2%
  \tagmcend%
  \tagstructend%
}
\newenvironment{examproblems}[2]{%
  \def\ProblemHeadingTag{#2}%
  \AssignTaggingSocketPlug{block/list/label}{examproblem}%
  \begin{enumerate}%
  \setcounter{enumi}{#1}%
  \renewcommand{\labelenumi}{\textbf{\theenumi.}}%
  \setlength{\topsep}{0.35em}%
  \setlength{\partopsep}{0pt}%
  \setlength{\itemsep}{0.48em}%
  \setlength{\parsep}{0.18em}%
}{\end{enumerate}}
\newenvironment{examsubparts}{%
  \AssignTaggingSocketPlug{block/list/label}{default}%
  \begin{enumerate}%
  \setlength{\topsep}{0.18em}%
  \setlength{\partopsep}{0pt}%
  \setlength{\itemsep}{0.16em}%
  \setlength{\parsep}{0.1em}%
}{\end{enumerate}}
\newenvironment{examinstructions}{%
  \par\begingroup\itshape
}{%
  \par\endgroup\vspace{0.18em}
}
\makeatletter
\renewcommand\subsection{\@startsection{subsection}{2}{0pt}%
  {-1.25ex plus -.25ex minus -.1ex}%
  {0.55ex plus .1ex}%
  {\normalfont\large\bfseries}}
\renewcommand{\@maketitle}{%
  \newpage\null\vskip -1em%
  {\footnotesize\bfseries
    \href{https://www.ufl.edu/}{University of Florida}, \href{https://math.ufl.edu/}{Department of Mathematics}\par}%
  \vskip -0.5em%
  \tagstructbegin{tag=H1,title={Analysis PhD Exam, August 2022}}%
  \tagpdfparaOff%
  \tagmcbegin{tag=H1}%
    {\LARGE\bfseries\href{https://gma.math.ufl.edu/exams/analysis-phd/}{Analysis PhD Exam}, August 2022\par}%
  \tagmcend%
  \tagpdfparaOn%
  \tagstructend%
  \vskip 0.25em%
  \tagmcbegin{artifact}%
    \hrule height 1pt%
  \tagmcend%
  \vskip 1em%
}
\makeatother
\title{Analysis PhD Exam, August 2022}
\author{}
\date{}
\begin{document}
\maketitle
\thispagestyle{empty}
\begin{examinstructions}
Write solutions in a neat and logical fashion, giving complete reasons for all steps. You must attempt SIX problems.
\end{examinstructions}
\begin{examproblems}{0}{H2}
\def\ProblemHeadingText{Problem 1.}
\item
Let \((X,\mathcal{M})\) be a measurable space and~\(\nu\) a signed measure. Define the Hahn decomposition of~\(X\) and the Jordan decomposition of~\(\nu\). Prove that the Jordan decomposition is unique.
\def\ProblemHeadingText{Problem 2.}
\item
Let \(\mu,\nu\) be positive measures defined on the same measure space. Define what it means for~\(\nu\) to be absolutely continuous with respect to~\(\mu\), and state the Radon–Nikodym theorem.
\def\ProblemHeadingText{Problem 3.}
\item
State Tonelli’s theorem, and give an example to show that the~\(\sigma\)-finiteness hypothesis is necessary.
\def\ProblemHeadingText{Problem 4.}
\item
Give an example of each of the following, if possible. (If not possible, give a brief explanation of why.) (When giving examples, be clear about what measure space you are working in.)
\begin{examsubparts}
\item[\textnormal{a)}]
a sequence of functions converging in measure but not in \(L^1\),
\item[\textnormal{b)}]
a sequence converging almost uniformly but not essentially uniformly,
\item[\textnormal{c)}]
a sequence converging in \(L^1\) but not almost everywhere.
\end{examsubparts}
\def\ProblemHeadingText{Problem 5.}
\item
Let \((X,\mathcal{M},\mu)\) be a~\(\sigma\)-finite measure space. Prove that the simple functions which belong to \(L^2(\mu)\) are dense in \(L^2(\mu)\).
\def\ProblemHeadingText{Problem 6.}
\item
Let~\(\mathcal{X}\) be a Banach space and \(\mathcal{N}\subset\mathcal{X}\) a closed subspace, and let \(\pi:\mathcal{X}\to\mathcal{X}/\mathcal{N}\) be the quotient map. Show that if \(f:\mathcal{X}\to\mathbb{F}\) is a bounded linear functional such that \(f(n)=0\) for all \(n\in\mathcal{N}\), then there is a bounded linear functional \(g:\mathcal{X}/\mathcal{N}\to\mathbb{F}\) such that \(f=g\circ\pi\).
\def\ProblemHeadingText{Problem 7.}
\item
Let~\(\mathcal{X}\) be a reflexive Banach space, with dual space \(\mathcal{X}^*\). Suppose \(B:\mathcal{X}\times\mathcal{X}^*\to\mathbb{F}\) is a bilinear mapping, and there is a constant \(C>0\) so that
\[
|B(x,x^*)|\le C\|x\|\|x^*\| \text{ for all } x\in\mathcal{X},\ x^*\in\mathcal{X}^*.
\]
 Prove that there is a bounded linear operator \(T:\mathcal{X}\to\mathcal{X}\) such that \(B(x,x^*)=x^*(Tx)\) for all \(x\in\mathcal{X},x^*\in\mathcal{X}^*\). Must~\(T\) be unique?
\def\ProblemHeadingText{Problem 8.}
\item
Let \(f:\mathbb{R}\to\mathbb{R}\) be a continuous function with compact support, and let \(\phi:\mathbb{R}\to\mathbb{R}\) be a nonnegative \(L^1\) function with \(\int_{-\infty}^{\infty}\phi(t)\,dt=1\). For each real \(\lambda>0\), consider the function
\[
f_\lambda(x):=\frac{1}{\lambda}\int_{-\infty}^{\infty}\phi\left(\frac{x-t}{\lambda}\right)f(t)\,dt.
\]
 Prove that \(f_\lambda\to f\) uniformly on~\(\mathbb{R}\) as \(\lambda\to0\).
\end{examproblems}
\end{document}
